\documentclass[10pt, aspectratio=169,t,xcolor=table]{beamer}

\setlength{\itemsep}{0pt}
\setlength{\parskip}{0pt}
\setlength{\parsep}{0pt}

\usetheme{Finalyear}

\usepackage{multicol}
\usepackage{pdfpages}
\usepackage{appendixnumberbeamer}
\usepackage{tabularx}
\usepackage{adjustbox}
\usepackage{enumitem}
% \setlist{nosep, leftmargin=*}
\newcommand{\shadowsq}{%
    \begin{tikzpicture}
        \fill[gray!90] (0.02,-0.05) rectangle (0.12,0.05);
        \fill[black] (-0.02,-0.02) rectangle (0.08,0.08);
        % \draw[white, very thin] (0,0.25) -- (0,0) -- (0.25,0);
    \end{tikzpicture}%
}
\setlist{label=\shadowsq, itemsep=0.1pt, parsep=1pt, topsep=5pt, leftmargin=8pt}
\usepackage{fontawesome5}
% \usepackage[
%     style=authoryear,
%     backend=biber,
%     language=english,
%     dashed=false
% ]{biblatex}
% \AtBeginBibliography{\footnotesize}
% \ExecuteBibliographyOptions{dashed=false}
% \addbibresource{references.bib}
% \setbeamertemplate{bibliography entry title}{}
% \setbeamertemplate{bibliography entry location}{}
% \setbeamertemplate{bibliography entry note}{}
% \setbeamercolor{bibliography item}{fg=blue!80}
% \setbeamercolor{bibliography entry author}{fg=black}
% \setbeamercolor{bibliography entry title}{fg=black}

% \setbeamertemplate{bibliography item}[triangle] % icons
% \setbeamercolor{cite}{fg=red!10!black}
% \DeclareCiteCommand{\cite}[\mkbibbrackets]
%   {\usebibmacro{prenote}}%
%   {\usebibmacro{citeindex}%
%     \printtext[bibhyperref]{\tiny \textbf{\printnames{labelname}, \printfield{year}}}}%
%   {\multicitedelim}%
%   {\usebibmacro{postnote}}%
% \setbeamerfont{frametitle}{size=\fontsize{11}{11}\selectfont}
% \setbeamercolor{bibliography item}{fg=black!80}
% \setbeamercolor{footnote mark}{fg=blue!80}

% Open refrence from local dir
% \newcommand{\pdflink}[2]{\href{../papers/#1.pdf}{#2}}
% \newcommand{\pdflink}[2]{\href{run:../papers/#1.pdf}{#2}}
% \newcommand{\pdflink}[2]{\href{file:papers/#1.pdf}{#2}}

\definecolor{darkgreen}{RGB}{0,100,0}
\definecolor{darkred}{RGB}{139,0,0}

% Theorem environments
\newtheorem{conj}{Conjecture}
\newtheorem{defi}{Definition}
\newtheorem{teo}{Theorem}

\setbeamertemplate{theorems}[numbered]
\setbeamertemplate{caption}[numbered]
\setcounter{figure}{0}

\setPrimaryColor{MyBlue}
\setLogos{lib/logo/iist-logo-horizontal.png}{lib/logo/iist-logo-square.png}

\begin{document}
% -------------------------------------------------------------------
%  TITLE INFORMATION
\title[MD Land-Atmosphere Coupling]{Land-Atmosphere Coupling in Monsoon Depressions via Soil Moisture Feedback}
\subtitle{Impact of Land-Surface Feedback on Post-Landfall Sustainability of Monsoon Depressions}
\date{}
% \author[Harshit P. D.]{
% 	\begin{columns}
% 		\column{0.4\textwidth}
% 		\centering
% 		\textbf{Presented by:} \\
% 		\textbf{Harshit Prashant Dhanwalkar} \\
% 		\tiny{SC21B164} \\
% 		\tiny{Earth System Sciences (ESS)}
%
% 		\column{0.4\textwidth}
% 		\centering
% 		\textbf{Supervisor:} \\
% 		Dr.\ A.\ Chandrasekar
% 	\end{columns}
% 	\hfill
% }
\setbeamertemplate{title page}{
	\vfill
	\begin{centering}
		\usebeamerfont{title}\usebeamercolor[fg]{title}\inserttitle\par
		\vspace{0.2cm}
		\usebeamerfont{subtitle}\usebeamercolor[fg]{subtitle}\insertsubtitle\par
		\vspace{0.5cm}

		\begin{columns}[T]
			\column{0.45\textwidth}
			\centering
			\textbf{Presented by:} \\
			\textbf{Harshit Prashant Dhanwalkar} \\
			\tiny{SC21B164} \\
			\tiny{Earth System Sciences (ESS)}

			\column{0.45\textwidth}
			\centering
			\textbf{Supervisor:} \\
			Dr.\ A.\ Chandrasekar
		\end{columns}
		\vfill
	\end{centering}
}
\frame[noframenumbering]{\titlepage}

% -------------------------------------------------------------------
%  TABLE OF CONTENTS
% \setLayout{plain}
% \begin{frame}[plain, c]{Table of Contents}
% 	\vspace{-0.3cm}
% 	\setbeamerfont{section in toc}{size=\fontsize{7}{8}\selectfont}
% 	\setbeamerfont{subsection in toc}{size=\fontsize{5.5}{6.5}\selectfont}
% 	\setbeamertemplate{section in toc}[sections numbered]
% 	\setbeamertemplate{subsection in toc}{\leavevmode\leftskip=0.6em\inserttocsubsection\par}
% 	\begin{multicols}{3}
% 		\tableofcontents
% 	\end{multicols}
% \end{frame}

% -------------------------------------------------------------------
\section{Introduction}
\setLayout{mainpoint}
\begin{frame}{}
	\frametitle{INTRODUCTION}
\end{frame}

% -------------------------------------------------------------------
\setLayout{vertical}
\begin{frame}{Monsoon Depression Inland Longevity}
	\footnotesize
	\begin{itemize}
		\item \textbf{Monsoon Depressions (MDs):} Large, slow-moving low-pressure systems
		      forming over the Bay of Bengal that are responsible for $\mathbf{>50\%}$ of
		      total Indian monsoon rainfall, with a typical lifespan of $\sim$2–7 days.
		      \vspace{0.25cm}
		\item \textbf{The Longevity Paradox:} Some MDs decay within 3 days of landfall.
		      Others persist for 7+ days deep inland - far from any oceanic moisture source.
		      What drives this difference?
		      \vspace{0.25cm}
		\item \textbf{The Brown Ocean Hypothesis:} Wet, clay-rich soils of the
		      Indo-Gangetic Plain mimic the ocean's latent heat supply - providing a
		      \textit{secondary energy source} that can sustain an MD long after it
		      leaves the Bay of Bengal.
		      \vspace{0.25cm}
		\item \textbf{Research Gap:} Despite well-studied atmospheric dynamics, the
		      \textit{quantitative role of land-atmosphere coupling} - specifically
		      soil moisture $\leftrightarrow$ precipitation feedback - remains poorly
		      constrained in operational forecasting.
	\end{itemize}
	\vspace{0.2cm}
	\begin{alertblock}{\small Practical Significance}
		\scriptsize
		Duration directly determines total rainfall, flood potential, and forecast skill.
		Current models have useful skill only to $\sim$72–96 h post-landfall.
		Land-surface factors are systematically underrepresented.
	\end{alertblock}
\end{frame}

% -------------------------------------------------------------------
\setLayout{horizontal}
\begin{frame}{The Brown Ocean Effect - Physical Mechanism}
	\begin{columns}
		\column{0.92\textwidth}
		% \vspace{-0.3cm}
		\begin{block}{\small What is the Brown Ocean Effect?}
			\footnotesize
			Over inland, soils act as a \textbf{``pseudo-ocean''} energy reservoir:
			\begin{itemize}
				\item Volumatric Soil Water Layer (SWVL1) $\rightarrow$ enhanced evapotranspiration (ET)
				\item ET $\rightarrow$ large latent heat flux (LHF) to atmosphere
				\item LHF $\rightarrow$ increased atmospheric instability (CAPE)
				\item CAPE $\rightarrow$ deep convection \& precipitation feedback
				\item Total Precipitation (TP) $\rightarrow$ recharges soil moisture
			\end{itemize}
		\end{block}

		\vspace{0.2cm}
		\begin{block}{\small The Feedback}
			\centering
			\footnotesize
			\begin{tikzpicture}[
				node distance=1.5cm,
				box/.style={rectangle, draw, rounded corners=2pt, minimum width=1.2cm, minimum height=0.6cm, align=center},
				arrow/.style={-{Stealth[scale=1.2]}, thick}
				]
				\node[box] (tp) {TP};
				\node[box, right=of tp] (swvl1) {SWVL1};
				\node[box, right=of swvl1] (lhf) {LHF};

				\draw[arrow] (tp) -- node[above, font=\tiny] {recharge} (swvl1);
				\draw[arrow] (swvl1) -- node[above, font=\tiny] {ET} (lhf);
				\draw[arrow] (lhf.south) -- ++(0,-0.5) coordinate (corner)
				node[below, font=\tiny] at ($(lhf.south)!0.5!(corner)$) {}
				(corner) -- (corner -| tp.south) node[below, font=\tiny] at ($(corner)!0.5!(corner -| tp.south)$) {convection}
				(corner -| tp.south) -- (tp.south);
			\end{tikzpicture}
		\end{block}

		% \column{0.48\textwidth}
		% \vspace{-0.3cm}
		% \begin{block}{\small Key Metrics Developed}
		% 	\footnotesize
		% 	\begin{itemize}
		% 		\item \textbf{FSI} (Feedback Strength Index): TP–SWVL1 coupling strength
		% 		\item \textbf{CE} (Coupling Efficiency): How efficiently soil moisture
		% 		      is converted to sustained atmospheric heating
		% 		\item \textbf{ACF} (Autocorrelation at 24h): Soil moisture memory -
		% 		      does yesterday's wet soil persist today?
		% 		\item \textbf{TSI} (Topographic Sensitivity Index): Terrain's role in
		% 		      modulating feedback
		% 		\item \textbf{ASM} (Antecedent Soil Moisture): Pre-landfall soil state
		% 	\end{itemize}
		% \end{block}
	\end{columns}
\end{frame}

% % -------------------------------------------------------------------
% \setLayout{vertical}
% \begin{frame}{Research Objectives}
% 	\textbf{\small Primary Objective:}
% 	\footnotesize Quantify how land-surface conditions - soil moisture, soil type,
% 	terrain, and antecedent moisture - control post-landfall Monsoon Depression longevity.
%
% 	\vspace{0.3cm}
% 	\textbf{\small Six Testable Hypotheses:}
% 	\footnotesize
% 	\begin{enumerate}
% 		\item \textbf{H1 - Coupling Strength:} FSI $>$ 0.4 predicts long-lived MDs
% 		      ($\geq$96 h); FSI $<$ 0.1 predicts short-lived systems.
% 		\item \textbf{H2 - Diurnal Coupling:} Land-phase diurnal SHF amplitude
% 		      ($\sim$56 W/m$^2$) is $4\times$ stronger than ocean phase ($\sim$14 W/m$^2$).
% 		\item \textbf{H3 - Antecedent Moisture:} ASM strongly predicts FSI
% 		      ($r \approx 0.92$); wetter pre-monsoon soil $\to$ stronger coupling.
% 		\item \textbf{H4 - Landfall Transition:} Heat flux signature at landfall
% 		      distinguishes Ocean-to-Land from Land-origin tracks.
% 		\item \textbf{H5 - Soil Type:} Clay soils (IGP) sustain feedback; sandy coastal
% 		      soils break the loop via rapid runoff.
% 		\item \textbf{H6 - ASM Enables but Doesn't Ensure:} High ASM is necessary
% 		      but insufficient - requires clay texture, FSI $>$ 0.4, and ACF $>$ 0.3.
% 	\end{enumerate}
%
% 	\vspace{0.1cm}
% 	\begin{alertblock}{\small Null Hypothesis (H0)}
% 		\scriptsize
% 		Long-lived MDs show no systematic difference in TP–SWVL1 coupling,
% 		diurnal SHF, or soil memory compared to short-lived systems.
% 	\end{alertblock}
% \end{frame}
% -------------------------------------------------------------------

\section{Methodology}
% -------------------------------------------------------------------
\setLayout{mainpoint}
\begin{frame}{}
	\frametitle{METHODOLOGY}
\end{frame}

% -------------------------------------------------------------------
\setLayout{horizontal}
\begin{frame}{Dataset Information}
	\begin{columns}
		\column{0.8\textwidth}
		\begin{block}{\footnotesize ERA5 (ECMWF Reanalysis v5) Single Level Reanalysis Dataset} % \tiny \cite{hersbach2020era5}}
			\footnotesize
			\begin{itemize}
				\item \textbf{Source:} Copernicus Climate Data Store (CDS), ECMWF
				\item \textbf{Temporal Coverage:} June–September 2024 (4 months, Indian monsoon season)
				\item \textbf{Spatial Domain :} $10$--$30^\circ$ N, $70$--$95^\circ$ E (Indian subcontinent + Bay of Bengal)
				      % FIX: the spatial (31 km globally) and temporal (hourly) resolutions of ERA5
				\item \textbf{Resolution:} 0.25$^\circ$ $\times$ 0.25$^\circ$
				      ($\approx$27.75 km latitude $\times$ 24–28 km longitude over India)
				\item \textbf{Temporal frequency:} Hourly (1-hourly tracking)
				\item \textbf{Format:} NetCDF4; land-sea mask (LSM) applied for
				      ocean masking in analysis
			\end{itemize}
		\end{block}

		\column{0.2\textwidth}
		\begin{figure}
			\centering
			\includegraphics[width=1\textwidth]{lib/logo/era5-logo.png}
		\end{figure}

		\vspace{0.2cm}
		\begin{figure}
			\centering
			\includegraphics[width=1\textwidth]{lib/logo/Copernicus vecto def Europe's eyes on Earth.png}
		\end{figure}
	\end{columns}
\end{frame}

% % -------------------------------------------------------------------
% NO need of data info, cut short number of slides
% \setLayout{vertical}
% \begin{frame}{Dataset Information}
% 	\vspace{-0.1cm}
% 	\begin{columns}[T]
% 		\column{0.44\textwidth}
% 		\begin{block}{\small ERA5 Dataset Variables :}
% 			\begin{itemize}\footnotesize
% 				\item Mean Sea Level Pressure (\texttt{MSL})
% 				\item Surface Pressure (\texttt{SP})
% 				\item 10-Meter U-Wind Component (\texttt{U10})
% 				\item 10-Meter V-Wind Component (\texttt{V10})
% 				\item Mean Surface Latent Heat Flux (\texttt{AVG\_SLHF})
% 				\item Mean surface Sensible Heat Flux (\texttt{AVG\_ISHF})
% 				\item Total Precipitation (\texttt{TP})
% 				\item Volumetric Soil Water Layer-1 (\texttt{SWVL1})
% 				\item Soil Temperature Level-1 (\texttt{STL1})
% 				\item 2-Meter Temperature (\texttt{T2M})
% 				\item 2-Meter Dew Point Temperature (\texttt{D2M})
% 			\end{itemize}
% 		\end{block}
%
% 		\column{0.55\textwidth}
% 		\begin{block}{\small Derived Variables :}
% 			\begin{itemize}\footnotesize
% 				\item \textbf{10-meter Wind Speed (\texttt{W10}) :}
% 				      \footnotesize
% 				      \begin{equation}
% 					      \texttt{W10} = \sqrt{(\texttt{U10})^2 + (\texttt{V10})^2}
% 				      \end{equation}
% 				\item \textbf{Relative Vorticity (\texttt{$\zeta$}):}
% 				      \footnotesize
% 				      \begin{equation}
% 					      \zeta = \frac{\partial(\texttt{V10)}}{\partial x} - \frac{\partial(\texttt{U10)}}{\partial y}
% 				      \end{equation}
% 				      % \item \textbf{Total Heat Flux (\texttt{THF}) :}
% 				      %       \footnotesize
% 				      %       \begin{equation}
% 				      % 	      \texttt{THF} = \texttt{SLHF} + \texttt{SSHF}
% 				      %       \end{equation}
% 				\item \textbf{Water Vapor Mixing Ratio (\texttt{w}) :}
% 				      \footnotesize
% 				      \begin{equation}
% 					      \texttt{w} = 0.622 \frac{e_s}{\texttt{SP} - e_s} \times 1000
% 				      \end{equation}
% 				      Where \textbf{Saturation Vapor Pressure (\texttt{$e_s$}) :}
% 				      \footnotesize
% 				      \begin{equation}
% 					      e_s = 6.112 \exp\left(\frac{17.67 * \texttt{D2M}}{\texttt{D2M} + 243.5}\right)
% 				      \end{equation}
% 				      and \texttt{D2M} = 2m Dew Point Temperature in $^\circ$C
% 			\end{itemize}
% 		\end{block}
% 	\end{columns}
% 	% \item \textbf{Processing :} Multi-month concatenation, temporal regularization, spatial clipping
% \end{frame}

% \begin{frame}{ERA5 Reanalysis Dataset Variables Selection}
% 	\begin{columns}
% 		\column{0.8\textwidth}
% 		\begin{block}{\small Key Variables}
% 			\footnotesize
% 			\begin{tabular}{p{0.30\linewidth} p{0.18\linewidth} p{0.40\linewidth}}
% 				\toprule
% 				\rowcolor{MyBlue!15}
% 				\textbf{Variable}        &                     & \textbf{Role}                   \\
% 				\midrule
% 				Mean sea-level pressure  & \texttt{msl}        & MD detection \& tracking        \\
% 				10m wind (u,v)           & \texttt{u10,v10}    & Circulation, speed              \\
% 				Soil water layer 1       & \texttt{swvl1}      & \textbf{Core feedback variable} \\
% 				Total precipitation      & \texttt{tp}         & \textbf{Main response variable} \\
% 				Latent heat flux         & \texttt{avg\_slhtf} & Moisture transport              \\
% 				Sensible heat flux       & \texttt{avg\_ishf}  & Surface heating                 \\
% 				2m temperature           & \texttt{t2m}        & Boundary layer                  \\
% 				2m dew point temperature & \texttt{d2m}        & Water vapor mixing ratio        \\
% 				Soil temperature L1      & \texttt{stl1}       & Land energy balance             \\
% 				\bottomrule
% 			\end{tabular}
% 		\end{block}
% 	\end{columns}
% \end{frame}

% -------------------------------------------------------------------
\subsection{MD Detection \& Tracking}
\setLayout{horizontal}
\begin{frame}{MD Detection \& Tracking Algorithm}
	\vspace{-0.4cm}
	\begin{columns}[T]
		\column{0.57\textwidth}
		\begin{block}{\small MD Detections Criteria :}
			\begin{itemize}\footnotesize
				\item \textbf{Pressure filter:} \texttt{MSL} local minimum $<$ 1000 hPa with a
				      $\Delta p \geq 2$ hPa depression relative to surroundings.
				\item \textbf{Contour Closure :} Closed isobar within \texttt{5$^{\circ}$} buffer
				\item \textbf{Track Duration :} \texttt{$\geq$ 24 hours}
				\item \textbf{Min net displacement:} \texttt{250 km}
				\item \textbf{Wind speed range\footnotemark[1] :} \texttt{5} to \texttt{19 m/s}
				\item \textbf{Total Precipitation (\texttt{TP}) :} \texttt{> 0.3 mm/hr} within buffer
				\item \textbf{Vorticity (\texttt{$\zeta$}) :} \texttt{> 0 /s} (cyclonic)
				\item \textbf{Max MD center translational speed :} \texttt{12 m/s}
				\item \textbf{Land-sea classification:} ERA5 Land–Sea Mask (\texttt{LSM}) applied to
				      tag each track point as ocean or land phase.
			\end{itemize}
		\end{block}

		\vfill
		\footnotetext[1]{\tiny Wind speed criteria based on IMD Monsoon Report 2024 (page 284), Source : \href{https://mausam.imd.gov.in/imd\_latest/monsoonreport2024.pdf}{\url{https://mausam.imd.gov.in/imd\_latest/monsoonreport2024.pdf}}}
		\footnotetext[2]{\tiny SWVL1 range based on ERA5 climatology for Indian region, Source : \href{https://www.researchgate.net/publication/365632311\_Assessment\_of\_ERA5-Land\_volumetric\_soil\_water\_layer\_product\_using\_in-situ\_and\_SMAP\_soil\_moisture\_observations}{\url{https://www.researchgate.net/publication/365632311\_Assessment\_of\_ERA5-Land\_volumetric\_soil\_water\_layer\_product\_using\_in-situ\_and\_SMAP\_soil\_moisture\_observations}}}

		\column{0.43\textwidth}
		\begin{block}{\small Tracking \& Post-Processing :}
			\footnotesize
			\begin{itemize}
				\item \textbf{Hourly Linking :} \texttt{50 km} search radius
				\item \textbf{Gap Bridging :} Up to \texttt{24 h}
				\item \textbf{Merge Overlaps :} \texttt{250 km} radius
				\item \textbf{Filtering :} WNW preference (\texttt{240$^{\circ}$ – 30$^{\circ}$})
			\end{itemize}
			\begin{itemize}
				\item \textbf{Landfall Detection :}
				      \begin{itemize}
					      \item \textbf{Database:} Natural Earth Physical Mask
					      \item \textbf{Resolution:} \texttt{50m} (1:50M scale)
					      \item \textbf{Method:} Point-in-polygon check
				      \end{itemize}
				\item \textbf{\texttt{SWVL1} Range\footnotemark[2] :} \texttt{0.02 - 0.55 m$^3$/m$^3$}
			\end{itemize}
		\end{block}

		\vspace{0.2cm}
		\begin{block}{\small Track Characterisation}
			\begin{itemize}\footnotesize
				\item \textbf{Translation speed:} Geodesic (Haversine) distance between
				      consecutive 1-hourly positions by $\Delta t$.
				\item \textbf{Landfall detection:} First ocean $\to$ land transition (diagonal
				      boundary algorithm; 0–25 km gap correction applied).
				\item \textbf{Temporal resolution:} 1-hourly vs.\ IMD 3-hourly - minor
				      start/end time differences are expected.
			\end{itemize}
		\end{block}
	\end{columns}
	%
	% 	% \vfill
	% 	% \footnotetext[1]{\tiny Wind speed criteria based on IMD Monsoon Report 2024 (page 284), Source : \href{https://mausam.imd.gov.in/imd\_latest/monsoonreport2024.pdf}{\url{https://mausam.imd.gov.in/imd\_latest/monsoonreport2024.pdf}}}
	% 	% \footnotetext[2]{\tiny SWVL1 range based on ERA5 climatology for Indian region, Source : \href{https://www.researchgate.net/publication/365632311\_Assessment\_of\_ERA5-Land\_volumetric\_soil\_water\_layer\_product\_using\_in-situ\_and\_SMAP\_soil\_moisture\_observations}{\url{https://www.researchgate.net/publication/365632311\_Assessment\_of\_ERA5-Land\_volumetric\_soil\_water\_layer\_product\_using\_in-situ\_and\_SMAP\_soil\_moisture\_observations}}}
\end{frame}

% -------------------------------------------------------------------
\setLayout{vertical}
\begin{frame}{The Feedback Mechanism}
	\textbf{\small Precipitation $\rightarrow$ Soil Moisture $\rightarrow$ Instability $\rightarrow$ Convection Loop}
	\vspace{0.2cm}
	\begin{block}{\small Diurnal Recharge Cycle (Land Phase)}
		\footnotesize
		\begin{enumerate}
			\item \textbf{Daytime (SHF peak):} Low thermal energy from moist soil drives SHF amplitude.
			      Warms boundary layer $\to$ maintains low-level convergence.
			\item \textbf{Nighttime (LHF shift):} Soil moisture reservoir sustains
			      evapotranspiration outside peak heating hours.
			      Prevents drying of the MD vortex core overnight.
		\end{enumerate}
		The self-sustaining loop (Brown Ocean Effect) allows MDs to persist
		for days over the Indian subcontinent \textit{without oceanic input}.
	\end{block}

	% FIX: Check elevation
	\begin{alertblock}{\small Elevation as Feedback Breaker}
		\scriptsize
		Eastern Ghats terrain ($>$900 m): orographic precipitation triggers rapid runoff
		on steep slopes ($>$8$^\circ$), preventing infiltration $\to$ soil stays dry
		$\to$ feedback loop broken $\to$ MD decays (Tracks 1, 4).
	\end{alertblock}
\end{frame}

% -------------------------------------------------------------------
\section{Results}
\setLayout{mainpoint}
\begin{frame}{}
	\frametitle{RESULTS}
\end{frame}

% -------------------------------------------------------------------
\subsection{Track Overview}
\setLayout{plain}
\begin{frame}{MD Track Map (June–September 2024)}
	\vspace{-0.5cm}
	\begin{columns}[T, onlytextwidth]
		\column{0.48\textwidth}
		\begin{figure}
			\centering
			\includegraphics[width=1.05\linewidth]{../../../version29-with-ocean-masking-lsm/output_2024/MD_track_2024.png}
		\end{figure}
		{\tiny \textbf{Fig :} ERA5-detected MD tracks for monsoon season 2024.}

		\column{0.48\textwidth}
		\begin{figure}
			\centering
			\includegraphics[width=1.35\linewidth]{../images/IMD-report-2024.png}
		\end{figure}
		{\tiny \textbf{Fig :} IMD Annual Monsoon Report 2024.}
	\end{columns}
\end{frame}

% -------------------------------------------------------------------
% TODO: Check
\setLayout{plain}
\begin{frame}{Validation - ERA5 detected MD track from IMD Annual Monsoon Report 2024}
	\vspace{-0.2cm}
	\begin{block}{\small Track Details}
		\footnotesize
		\begin{tabular}{p{0.1\linewidth}p{0.2\linewidth}p{0.3\linewidth}p{0.2\linewidth}}
			\toprule
			\rowcolor{MyBlue!15}
			\textbf{Track} & \textbf{Period} & \textbf{Origin} & \textbf{Duration} ($\geq$96h/Long) \\
			\midrule
			1              & 18–21 Jul       & Ocean-Origin    & 78h / Short                        \\
			2              & 01–06 Aug       & Land-Origin     & 106h / Long                        \\
			3              & 23–29 Aug       & Land-Origin     & 146h / Long                        \\
			4              & 30 Aug–02 Sep   & Ocean-Origin    & 75h / Short                        \\
			5              & 06–13 Sep       & Ocean-Origin    & 165h / Long                        \\
			6              & 13–18 Sep       & Land-Origin     & 117h / Long                        \\
			\bottomrule
		\end{tabular}
	\end{block}

	\vspace{0.2cm}
	\begin{alertblock}{\small Validation Result}
		\footnotesize
		Minor $\pm$6–12h differences in start/end times arise from
		ERA5 hourly vs.\ IMD 3-hourly resolution.
		Mean detected duration: $\mathbf{115.2 \pm 38.3}$ hours.
	\end{alertblock}
\end{frame}

% -------------------------------------------------------------------
% --- TRACK 1 -------------------------------------------------
\subsection{Track 1 Short-lived, Ocean-Origin, July}
\setLayout{plain}
\begin{frame}{Track 1 - Short-lived, Ocean-Origin, July}
	\vspace{-0.3cm}
	\begin{center}
		\begin{columns}[T]
			\column{0.5\textwidth}
			\includegraphics[width=0.85\textwidth]{../../../version29-with-ocean-masking-lsm/output_2024/track_1/MD_track_1_detailed.png}

			\begin{block}{\small Statistics}
				\vspace{-0.05cm}
				\begin{itemize}\footnotesize
					\item \textbf{Period:} 18 Jul 04:00 – 21 Jul 07:00, 2024
					\item \textbf{Duration} : 78 hours ($\sim$ 3 days)
					\item \textbf{Category} : short-lived MD
					\item \textbf{Origin} : Bay of Bengal (Oceanic)
					\item \textbf{Path:} Ocean $\to$ Eastern Ghats / Odisha coast
				\end{itemize}
			\end{block}

			\column{0.5\textwidth}
			\includegraphics[width=0.75\textwidth]{../images/July-19-20-2024.png}

			{\tiny RSMC New Delhi Report 2024}\footnote{\tiny Source RSMC New Delhi Report 2024 : \href{https://rsmcnewdelhi.imd.gov.in/report.php?internal\_menu=MzQ=}{\url{https://rsmcnewdelhi.imd.gov.in/report.php?internal\_menu=MzQ=}}}
		\end{columns}
	\end{center}
\end{frame}

% -------------------------------------------------------------------
\setLayout{plain}
\begin{frame}{Track 1 Heat Flux \&  Precipitation \& Soil Moisture Time Evolution}
	\begin{columns}[T]
		\column{0.50\textwidth}
		\vspace{-0.3cm}
		\begin{figure}
			\centering
			\includegraphics[width=1.1\textwidth]{../../../version29-with-ocean-masking-lsm/output_2024/track_1/time_series/track_1_heat_flux_time_series.png}
		\end{figure}
		{\tiny \textbf{Fig :} Heat flux evolution, Track 1. Clear diurnal SHF oscillations throughout land phase confirm active boundary layer coupling and not due to ocean-phase suppression.}

		\column{0.50\textwidth}
		\vspace{-0.3cm}
		\begin{figure}
			\centering
			\includegraphics[width=1.1\textwidth]{../../version25-with-ocean-masking-lsm/output_2024/track_1/time_series/track_1_TP-SWVL1_time_series.png}
		\end{figure}
		{\tiny \textbf{Fig :} TP and SWVL1 for Track 1.
		High baseline SWVL1 ($\sim$0.4 m$^3$/m$^3$) supports sustained feedback even as individual TP pulses fluctuate.}
	\end{columns}
\end{frame}

% -------------------------------------------------------------------
\subsubsection{Track 1 Relation}
\setLayout{plain}
\begin{frame}{Track 1 Soil Moisture Heat flux, Precipitation Soil Moisture \& Precipitation Heat Flux Relation}
	\vspace{-0.3cm}
	\begin{columns}[T]
		\begin{column}{0.5\textwidth}
			\begin{center}
				\includegraphics[width=1.1\textwidth]{../../version29-with-ocean-masking-lsm/output_2024/track_1/correlation_analysis/relation/track_1_SWVL1-HF_relation.png} \\
				\includegraphics[width=1.1\textwidth]{../../version29-with-ocean-masking-lsm/output_2024/track_1/correlation_analysis/relation/track_1_TP-HF_relation.png}
			\end{center}
		\end{column}

		\begin{column}{0.5\textwidth}
			\begin{center}
				\includegraphics[width=1.1\textwidth]{../../version29-with-ocean-masking-lsm/output_2024/track_1/correlation_analysis/relation/track_1_TP-SWVL1_relation.png}
			\end{center}
		\end{column}
	\end{columns}
\end{frame}

% -------------------------------------------------------------------
\subsubsection{Track 1 Hovm\"oller}
\includepdf[
	pages=-,
	scale=1,
	frame=true,
	pagecommand={\stepcounter{framenumber}},
	offset=0mm 0mm
]{../slides/hov-1.pdf}

% -------------------------------------------------------------------
\subsubsection{Track 1 SRTM Elevation Map}
\setLayout{plain}
\begin{frame}{Track 1 SRTM Elevation Map}
	\vspace{-0.3cm}
	\begin{center}
		\includegraphics[width=0.6\textwidth]{../../version29-with-ocean-masking-lsm/output_2024/track_1/elevation_analysis/MD_track_1_elevation_mosaic.png}
	\end{center}
\end{frame}

% -------------------------------------------------------------------
\subsubsection{Track 1 Translation Speed}
\setLayout{horizontal}
\begin{frame}{Track 1 - Translation Speed of MD Center}
	\vspace{-0.3cm}
	\begin{center}
		\includegraphics[width=0.8\textwidth]{../../../version29-with-ocean-masking-lsm/output_2024/track_1/translational-speed/track_1_speed_ms.png}

		\begin{block}{\footnotesize Interpretation}
			\scriptsize As Track makes landfall over the Odisha coast, the MD center translation speed drops due surface friction and reduced pressure gradient over land. The inland penetration shows a slight recovery in speed before the system dissipates, reflecting the absence of any sustained land-surface feedback to anchor the vortex.
		\end{block}
	\end{center}
\end{frame}

% --- TRACK 2 -------------------------------------------------
\subsection{Track 2 Long-lived, Land-Origin, August}
\setLayout{plain}
\begin{frame}{Track 2 - Long-lived, Land-Origin, August}
	\vspace{-0.3cm}
	\begin{center}
		\begin{columns}[T]
			\column{0.5\textwidth}
			\includegraphics[width=0.95\textwidth]{../../../version29-with-ocean-masking-lsm/output_2024/track_2/MD_track_2_detailed.png}

			\begin{block}{\small Statistics}
				\vspace{-0.05cm}
				\begin{itemize}\footnotesize
					\item \textbf{Period:} 01 Aug 19:00 – 06 Aug 04:00, 2024
					\item \textbf{Duration:} 106 h ($\sim$4.5 days)
					\item \textbf{Category} : Long-lived MD
					\item \textbf{Origin:} Over land (Madhya Pradesh / Deep Depression)
				\end{itemize}
			\end{block}

			\column{0.5\textwidth}
			\includegraphics[width=0.75\textwidth]{../images/August-02-06-2024.png}

			{\tiny RSMC New Delhi Report 2024}% \footnote{\tiny Source RSMC New Delhi Report 2024 : \href{https://rsmcnewdelhi.imd.gov.in/report.php?internal\_menu=MzQ=}{\url{https://rsmcnewdelhi.imd.gov.in/report.php?internal\_menu=MzQ=}}}
		\end{columns}

		\begin{block}{\footnotesize Interpretation}
			\scriptsize Madhya Pradesh clay soils retain moisture efficiently $\to$ sustained LHF $\to$ long inland survival.
			The ``Deep Depression'' formed far from coast is evidence of land-surface flux acting as secondary energy source.
		\end{block}
	\end{center}
\end{frame}

% -------------------------------------------------------------------
\setLayout{plain}
\begin{frame}{Track 2 Heat Flux \&  Precipitation \& Soil Moisture Time Evolution}
	\begin{columns}[T]
		\column{0.50\textwidth}
		\vspace{-0.3cm}
		\begin{figure}
			\centering
			\includegraphics[width=1.1\textwidth]{../../../version29-with-ocean-masking-lsm/output_2024/track_2/time_series/track_2_heat_flux_time_series.png}
		\end{figure}
		{\tiny \textbf{Fig :} Heat flux time series. Land-origin track shows sustained positive SHF (diurnal pumping) throughout - no ocean-phase suppression.}

		\column{0.50\textwidth}
		\vspace{-0.3cm}
		\begin{figure}
			\centering
			\includegraphics[width=1.1\textwidth]{../../version25-with-ocean-masking-lsm/output_2024/track_2/time_series/track_2_TP-SWVL1_time_series.png}
		\end{figure}
		{\tiny \textbf{Fig :} TP and SWVL1 time series. High baseline SWVL1 ($\sim$0.45 m$^3$/m$^3$) supports sustained feedback even as individual TP pulses fluctuate.}
	\end{columns}
\end{frame}

% -------------------------------------------------------------------
\subsubsection{Track 2 Relation}
\setLayout{plain}
\begin{frame}{Track 2 Soil Moisture Heat flux, Precipitation Soil Moisture \& Precipitation Heat Flux Relation}
	\vspace{-0.3cm}
	\begin{columns}[T]
		\begin{column}{0.5\textwidth}
			\begin{center}
				\includegraphics[width=1.1\textwidth]{../../version29-with-ocean-masking-lsm/output_2024/track_2/correlation_analysis/relation/track_2_SWVL1-HF_relation.png} \\
				\includegraphics[width=1.1\textwidth]{../../version29-with-ocean-masking-lsm/output_2024/track_2/correlation_analysis/relation/track_2_TP-HF_relation.png}
			\end{center}
		\end{column}

		\begin{column}{0.5\textwidth}
			\begin{center}
				\includegraphics[width=1.1\textwidth]{../../version29-with-ocean-masking-lsm/output_2024/track_2/correlation_analysis/relation/track_2_TP-SWVL1_relation.png}
			\end{center}
		\end{column}
	\end{columns}
\end{frame}

% -------------------------------------------------------------------
\subsubsection{Track 2 Hovm\"oller}
\includepdf[
	pages=-,
	scale=1,
	frame=true,
	pagecommand={\stepcounter{framenumber}},
	offset=0mm 0mm
]{../slides/hov-2.pdf}

% -------------------------------------------------------------------
\subsubsection{Track 2 SRTM Elevation Map}
\setLayout{plain}
\begin{frame}{Track 2 SRTM Elevation Map}
	\vspace{-0.3cm}
	\begin{center}
		\includegraphics[width=0.6\textwidth]{../../version29-with-ocean-masking-lsm/output_2024/track_2/elevation_analysis/MD_track_2_elevation_mosaic.png}
	\end{center}
\end{frame}

% -------------------------------------------------------------------
\subsubsection{Track 2 Translation Speed}
\begin{frame}{Track 2 - Translation Speed of MD Center}
	\vspace{-0.3cm}
	\begin{center}
		\includegraphics[width=0.8\textwidth]{../../../version29-with-ocean-masking-lsm/output_2024/track_2/translational-speed/track_2_speed_ms.png}

		\begin{block}{\footnotesize Interpretation}
			\scriptsize Track originates over land and shows a deceleration phase as the system organises over Central India. Translation speed subsequently increases as the MD moves further inland, consistent with steering by the large-scale monsoon flow. The sustained slow-to-moderate speed over the Indo-Gangetic Plain allows prolonged soil moisture recharge, contributing to its 106-hour lifespan.
		\end{block}
	\end{center}
\end{frame}

% --- TRACK 3 -------------------------------------------------
\subsection{Track 3 Long-lived, Land-Origin, August}
\setLayout{plain}
\begin{frame}{Track 3 - Long-lived, Land-Origin, August}
	\vspace{-0.3cm}
	\begin{center}
		\begin{columns}[T]
			\column{0.5\textwidth}
			\includegraphics[width=0.9\textwidth]{../../../version29-with-ocean-masking-lsm/output_2024/track_3/MD_track_3_detailed.png}

			\begin{block}{\small Statistics}
				\vspace{-0.05cm}
				\begin{itemize}\footnotesize
					\item \textbf{Period:} 23 Aug 08:00 – 29 Aug 10:00, 2024
					\item \textbf{Duration:} 146 h ($\sim$6 days)
					\item \textbf{Category} : Long-lived MD
					\item \textbf{Origin:} Over land (Central India / Deep Depression)
				\end{itemize}
			\end{block}

			\column{0.5\textwidth}
			% no RSMC image
		\end{columns}

		% \vspace{0.2cm}
		% \begin{alertblock}{\small Key Finding}
		% 	\scriptsize
		% 	Best combination of high moisture capacity + moderate memory + strongest coupling $\to$ highest CE among all tracks. Validates H1 and H3: wet soil $\to$ strong FSI $\to$ long survival. 
		% \end{alertblock}
	\end{center}
\end{frame}

% -------------------------------------------------------------------
\setLayout{plain}
\begin{frame}{Track 3 Heat Flux \&  Precipitation \& Soil Moisture Time Evolution}
	\begin{columns}[T]
		\column{0.50\textwidth}
		\vspace{-0.3cm}
		\begin{figure}
			\centering
			\includegraphics[width=1.1\textwidth]{../../../version29-with-ocean-masking-lsm/output_2024/track_3/time_series/track_3_heat_flux_time_series.png}
		\end{figure}
		{\tiny \textbf{Fig :} Heat flux evolution. Clear diurnal SHF oscillations throughout land phase confirm active boundary layer coupling and not due to ocean-phase suppression.}

		\column{0.50\textwidth}
		\vspace{-0.3cm}
		\begin{figure}
			\centering
			\includegraphics[width=1.1\textwidth]{../../version25-with-ocean-masking-lsm/output_2024/track_3/time_series/track_3_TP-SWVL1_time_series.png}
		\end{figure}
		{\tiny \textbf{Fig :} TP and SWVL1 time series. High baseline SWVL1 ($\sim$0.4 m$^3$/m$^3$) supports sustained feedback even as individual TP pulses fluctuate.}
	\end{columns}
\end{frame}

% -------------------------------------------------------------------
\subsubsection{Track 3 Relation}
\setLayout{plain}
\begin{frame}{Track 3 Soil Moisture Heat Flux,  Precipitation Soil Moisture \& Precipitation Heat Flux Relation}
	\vspace{-0.3cm}
	\begin{columns}[T]
		\begin{column}{0.5\textwidth}
			\begin{center}
				\includegraphics[width=1.1\textwidth]{../../version29-with-ocean-masking-lsm/output_2024/track_3/correlation_analysis/relation/track_3_SWVL1-HF_relation.png} \\
				\includegraphics[width=1.1\textwidth]{../../version29-with-ocean-masking-lsm/output_2024/track_3/correlation_analysis/relation/track_3_TP-HF_relation.png}
			\end{center}
		\end{column}

		\begin{column}{0.5\textwidth}
			\begin{center}
				\includegraphics[width=1.1\textwidth]{../../version29-with-ocean-masking-lsm/output_2024/track_3/correlation_analysis/relation/track_3_TP-SWVL1_relation.png}
			\end{center}
		\end{column}
	\end{columns}
\end{frame}

% -------------------------------------------------------------------
\subsubsection{Track 3 Hovm\"oller}
\includepdf[
	pages=-,
	scale=1,
	frame=true,
	pagecommand={\stepcounter{framenumber}},
	offset=0mm 0mm
]{../slides/hov-3.pdf}

% -------------------------------------------------------------------
\subsubsection{Track 3 SRTM Elevation Map}
\setLayout{plain}
\begin{frame}{Track 3 SRTM Elevation Map}
	\vspace{-0.3cm}
	\begin{center}
		\includegraphics[width=0.6\textwidth]{../../version29-with-ocean-masking-lsm/output_2024/track_3/elevation_analysis/MD_track_3_elevation_mosaic.png}
	\end{center}
\end{frame}

% -------------------------------------------------------------------
\subsubsection{Track 3 Translation Speed}
\begin{frame}{Track 3 Translation Speed of MD Center}
	\vspace{-0.3cm}
	\begin{center}
		\includegraphics[width=0.8\textwidth]{../../../version29-with-ocean-masking-lsm/output_2024/track_3/translational-speed/track_3_speed_ms.png}

		\begin{block}{\footnotesize Interpretation}
			\scriptsize Track exhibits the most pronounced deceleration among all tracks, stalling near Central India before gradually accelerating inland. This near-stationary phase maximises rainfall accumulation over a single region, driving the strongest soil moisture recharge observed across all tracks. The subsequent speed increase corresponds to the system being picked up by the mid-tropospheric westerly flow.
		\end{block}
	\end{center}
\end{frame}

% --- TRACK 4 -------------------------------------------------
\subsection{Track 4 Short-lived, Ocean-Origin / Pre-Cyclone ASNA, August-September}
\setLayout{plain}
\begin{frame}{Track 4 - Short-lived, Ocean-Origin, Pre-Cyclone ASNA, Aug–Sep}
	\vspace{-0.45cm}
	\begin{center}
		\begin{columns}[T]
			\column{0.5\textwidth}
			\includegraphics[width=0.82\textwidth]{../../version29-with-ocean-masking-lsm/output_2024/track_4/MD_track_4_detailed.png}

			\begin{block}{\small Statistics}
				\vspace{-0.05cm}
				\begin{itemize}\footnotesize
					\item \textbf{Period:} 30 Aug 05:00 – 02 Sep 08:00, 2024
					\item \textbf{Duration:} 75 h ($\sim$3 days)
					\item \textbf{Category} : Short--lived MD
					\item \textbf{Origin:} Bay of Bengal (Oceanic)
					\item \textbf{Path:} Ocean $\to$ Eastern Ghats
				\end{itemize}
			\end{block}

			\column{0.5\textwidth}
			\includegraphics[width=0.8\textwidth]{../images/August-31-Sep-02-2024.png}

			{\tiny RSMC New Delhi Report 2024}% \footnote{\tiny Source RSMC New Delhi Report 2024 : \href{https://rsmcnewdelhi.imd.gov.in/report.php?internal\_menu=MzQ=}{\url{https://rsmcnewdelhi.imd.gov.in/report.php?internal\_menu=MzQ=}}}
		\end{columns}

		\begin{alertblock}{\footnotesize The Control Case}
			\scriptsize Despite higher intensity (Cyclone status), longevity was limited by absence of inland soil moisture feedback. Rapid coastal runoff (sandy soils) and high elevation prevents feedback loop.
		\end{alertblock}
	\end{center}
\end{frame}

% -------------------------------------------------------------------
\setLayout{plain}
\begin{frame}{Track 4 Heat Flux \&  Precipitation \& Soil Moisture Time Evolution}
	\begin{columns}[T]
		\column{0.50\textwidth}
		\vspace{-0.3cm}
		\begin{figure}
			\centering
			\includegraphics[width=1.1\textwidth]{../../../version29-with-ocean-masking-lsm/output_2024/track_4/time_series/track_4_heat_flux_time_series.png}
		\end{figure}
		{\tiny \textbf{Fig :} Heat flux evolution, Diurnal SHF oscillations over land phase}

		\column{0.50\textwidth}
		\vspace{-0.3cm}
		\begin{figure}
			\centering
			\includegraphics[width=1.1\textwidth]{../../version25-with-ocean-masking-lsm/output_2024/track_4/time_series/track_4_TP-SWVL1_time_series.png}
		\end{figure}
		{\tiny \textbf{Fig :} TP and SWVL1 time series
		TP pulse after landfall translate to sustained SWVL1 increase.}
	\end{columns}
\end{frame}

% -------------------------------------------------------------------
\subsubsection{Track 4 Relation}
\setLayout{plain}
\begin{frame}{Track 4 Soil Moisture Heat Flux, Precipitation Soil Moisture \& Precipitation Heat Flux Relation}
	\vspace{-0.3cm}
	\begin{columns}[T]
		\begin{column}{0.5\textwidth}
			\begin{center}
				\includegraphics[width=1.1\textwidth]{../../version29-with-ocean-masking-lsm/output_2024/track_4/correlation_analysis/relation/track_4_SWVL1-HF_relation.png} \\
				\includegraphics[width=1.1\textwidth]{../../version29-with-ocean-masking-lsm/output_2024/track_4/correlation_analysis/relation/track_4_TP-HF_relation.png}
			\end{center}
		\end{column}

		\begin{column}{0.5\textwidth}
			\begin{center}
				\includegraphics[width=1.1\textwidth]{../../version29-with-ocean-masking-lsm/output_2024/track_4/correlation_analysis/relation/track_4_TP-SWVL1_relation.png}
			\end{center}
		\end{column}
	\end{columns}
\end{frame}

% -------------------------------------------------------------------
\subsubsection{Track 4 Hovm\"oller}
\includepdf[
	pages=-,
	scale=1,
	frame=true,
	pagecommand={\stepcounter{framenumber}},
	offset=0mm 0mm
]{../slides/hov-4.pdf}

% -------------------------------------------------------------------
\subsubsection{Track 4 SRTM Elevation Map}
\setLayout{plain}
\begin{frame}{Track 4 SRTM Elevation Map}
	\vspace{-0.4cm}
	\begin{center}
		\includegraphics[width=0.6\textwidth]{../../version29-with-ocean-masking-lsm/output_2024/track_4/elevation_analysis/MD_track_4_elevation_mosaic.png}
	\end{center}
\end{frame}

% -------------------------------------------------------------------
\subsubsection{Track 4 Translation Speed}
\begin{frame}{Track 4 Translation Speed of MD Center}
	\vspace{-0.3cm}
	\begin{center}
		\includegraphics[width=0.8\textwidth]{../../../version29-with-ocean-masking-lsm/output_2024/track_4/translational-speed/track_4_speed_ms.png}

		\begin{block}{\footnotesize Interpretation}
			\scriptsize After landfall over the Odisha–Andhra coast, Track decelerates as it encounters the elevated Eastern Ghats terrain. The system then accelerates westward but finds no sustained soil moisture feedback over the sandy coastal soils and elevated Eastern Ghats. The combination of rapid translation and weak land-surface coupling leads to quick dissipation within 75 hours of formation.
		\end{block}
	\end{center}
\end{frame}

% --- TRACK 5 -------------------------------------------------
\subsection{Track 5 Long-lived, Ocean-origin, September}
\setLayout{plain}
\begin{frame}{Track 5 - Long-lived, Ocean-origin, September}
	\vspace{-0.3cm}
	\begin{center}
		\begin{columns}[T]
			\column{0.5\textwidth}
			\includegraphics[width=0.82\textwidth]{../../version29-with-ocean-masking-lsm/output_2024/track_5/MD_track_5_detailed.png}

			\begin{block}{\small Statistics}
				\vspace{-0.05cm}
				\begin{itemize}\footnotesize
					\item \textbf{Period:} 06 Sep 06:00 – 13 Sep 03:00, 2024
					\item \textbf{Duration:} 165 h ($\sim$7 days)
					\item \textbf{Category} : Long-lived md
					\item \textbf{Origin:} Bay Of Bengal (oceanic)
				\end{itemize}
			\end{block}

			\column{0.5\textwidth}
			\includegraphics[width=0.75\textwidth]{../images/september-08-10-2024.png}

			{\tiny RSMC New delhi report 2024}% \footnote{\tiny source rsmc new delhi report 2024 : \href{https://rsmcnewdelhi.imd.gov.in/report.php?internal\_menu=mzq=}{\url{https://rsmcnewdelhi.imd.gov.in/report.php?internal\_menu=mzq=}}}
		\end{columns}

		\begin{alertblock}{\footnotesize The anomaly - track 5 survived longest}
			\scriptsize Soil moisture persistence day-to-day sustains the vortex even during weakened coupling.
			Eastern Ghats slowdown (Sep 10) $\to$ vorticity recovers (Sep 11 IMD) $\to$ system re-energises from accumulated soil moisture over plains.
		\end{alertblock}
	\end{center}
\end{frame}

% -------------------------------------------------------------------
\setLayout{plain}
\begin{frame}{Track 5 Heat Flux, Precipitation \& Soil Moisture Time Evolution}
	\vspace{-0.4cm}
	\begin{columns}[T]
		\column{0.50\textwidth}
		\vspace{-0.4cm}
		\begin{figure}
			\centering
			\includegraphics[width=1.1\textwidth]{../../../version29-with-ocean-masking-lsm/output_2024/track_5/time_series/track_5_heat_flux_time_series.png}
		\end{figure}
		{\tiny \textbf{Fig :} Heat flux evolution. Clear diurnal SHF oscillations throughout land phase confirm active boundary layer coupling and not due to ocean-phase suppression. SHF turns positive over land - the surface actively heats the atmosphere.}

		\column{0.50\textwidth}
		\vspace{-0.4cm}
		\begin{figure}
			\centering
			\includegraphics[width=1.1\textwidth]{../../version25-with-ocean-masking-lsm/output_2024/track_5/time_series/track_5_TP-SWVL1_time_series.png}
		\end{figure}
		{\tiny \textbf{Fig :} TP and SWVL1 time series.	TP pulse after landfall translate to sustained SWVL1 increase.}
	\end{columns}
\end{frame}

% -------------------------------------------------------------------
\subsubsection{Track 5 Relation}
\setLayout{plain}
\begin{frame}{Track 5 Soil Moisture Heat Flux, Precipitation Soil Moisture \& Precipitation Heat Flux Relation}
	\vspace{-0.4cm}
	\begin{columns}[T]
		\column{0.5\textwidth}
		\vspace{-0.4cm}
		\begin{center}
			\includegraphics[width=1.1\textwidth]{../../version29-with-ocean-masking-lsm/output_2024/track_5/correlation_analysis/relation/track_5_SWVL1-HF_relation.png} \\
			\includegraphics[width=1.1\textwidth]{../../version29-with-ocean-masking-lsm/output_2024/track_5/correlation_analysis/relation/track_5_TP-HF_relation.png}
		\end{center}

		\column{0.5\textwidth}
		\vspace{-0.4cm}
		\begin{center}
			\includegraphics[width=1.1\textwidth]{../../version29-with-ocean-masking-lsm/output_2024/track_5/correlation_analysis/relation/track_5_TP-SWVL1_relation.png}
		\end{center}
	\end{columns}
\end{frame}

% -------------------------------------------------------------------
\subsubsection{Track 5 Hovm\"oller}
\includepdf[
	pages=-,
	scale=1,
	frame=true,
	pagecommand={\stepcounter{framenumber}},
	offset=0mm 0mm
]{../slides/hov-5.pdf}

% -------------------------------------------------------------------
\subsubsection{Track 5 SRTM Elevation Map}
\setLayout{plain}
\begin{frame}{Track 5 SRTM Elevation Map}
	\vspace{-0.4cm}
	\begin{center}
		\includegraphics[width=0.6\textwidth]{../../version29-with-ocean-masking-lsm/output_2024/track_5/elevation_analysis/MD_track_5_elevation_mosaic.png}
	\end{center}
\end{frame}

% -------------------------------------------------------------------
\subsubsection{Track 5 Translation Speed}
\begin{frame}{Track 5 Translation Speed of MD Center}
	\vspace{-0.3cm}
	\begin{center}
		\includegraphics[width=0.8\textwidth]{../../../version29-with-ocean-masking-lsm/output_2024/track_5/translational-speed/track_5_speed_ms.png}

		\begin{block}{\footnotesize Interpretation}
			\scriptsize Track shows a two-phase speed profile. Translation speed decreases as the system encounters the Eastern Ghats around September 10, effectively stalling the MD. This slowdown acts as a concentration zone - accumulated rainfall over elevated terrain recharges soil moisture on the plains under. Vorticity subsequently recovers (IMD, Sep 11) and the system accelerates inland, sustained by the soil moisture built up during the stall period. This pattern is the signature of Track 5's anomalous 165-hour longest survival.
		\end{block}
	\end{center}
\end{frame}

% --- TRACK 6 -------------------------------------------------
\subsection{Track 6 Long-lived, Land-Origin, September}
\setLayout{plain}
\begin{frame}{Track 6 - Long-lived, Land-Origin, September}
	\vspace{-0.3cm}
	\begin{center}
		\begin{columns}[T]
			\column{0.5\textwidth}
			\includegraphics[width=0.8\textwidth]{../../version29-with-ocean-masking-lsm/output_2024/track_6/md_track_6_detailed.png}

			\begin{block}{\small Statistics}
				\vspace{-0.05cm}
				\begin{itemize}\footnotesize
					\item \textbf{Period:} 13 Sep 12:00 – 18 Sep 07:00, 2024
					\item \textbf{Duration:} 117 h ($\sim$5 days)
					\item \textbf{Category:} Long-lived MD
					\item \textbf{Origin:} Over land (Eastern India / Deep Depression)
				\end{itemize}
			\end{block}

			\column{0.5\textwidth}
			\includegraphics[width=0.75\textwidth]{../images/September-13-18-2024.png}

			{\tiny RSMC New Delhi report 2024}% \footnote{\tiny source rsmc new delhi report 2024 : \href{https://rsmcnewdelhi.imd.gov.in/report.php?internal\_menu=mzq=}{\url{https://rsmcnewdelhi.imd.gov.in/report.php?internal\_menu=mzq=}}}
		\end{columns}

		% \begin{block}{\small Interpretation}
		% 	\scriptsize
		% 	Third consecutive land-origin Deep Depression over Central India.
		% 	Reinforces pattern: Tracks 2, 5, 6 over Madhya Pradesh exhibit
		% 	consistent feedback loops, intensifying into ``Deep'' systems
		% 	far from coast. Saturated soil (ACF near 0) means minimal
		% 	additional memory needed - soil is already at capacity.
		% \end{block}
	\end{center}
\end{frame}

% -------------------------------------------------------------------
\setLayout{plain}
\begin{frame}{Track 6 Heat Flux \&  Precipitation \& Soil Moisture Time Evolution}
	\begin{columns}[T]
		\column{0.50\textwidth}
		\vspace{-0.4cm}
		\begin{figure}
			\centering
			\includegraphics[width=1.1\textwidth]{../../../version29-with-ocean-masking-lsm/output_2024/track_6/time_series/track_6_heat_flux_time_series.png}
		\end{figure}
		{\tiny \textbf{Fig :} Heat flux evolution, Clear diurnal SHF oscillations throughout land phase confirm active boundary layer coupling and not due to ocean-phase suppression.}

		\column{0.50\textwidth}
		\vspace{-0.4cm}
		\begin{figure}
			\centering
			\includegraphics[width=1.1\textwidth]{../../version25-with-ocean-masking-lsm/output_2024/track_6/time_series/track_6_TP-SWVL1_time_series.png}
		\end{figure}
		{\tiny \textbf{Fig :} TP and SWVL1 time series. High baseline SWVL1 ($\sim$0.4 m$^3$/m$^3$) supports sustained feedback even as individual TP pulses fluctuate. TP pulse after landfall translate to sustained SWVL1 increase.}
	\end{columns}
\end{frame}

% -------------------------------------------------------------------
\subsubsection{Track 6 Relation}
\setLayout{plain}
\begin{frame}{Track 6 Soil Moisture heat flux,  Precipitation Soil Moisture \& Precipitation Heat Flux Relation}
	\begin{columns}[T]
		\column{0.50\textwidth}
		\vspace{-0.4cm}
		\begin{center}
			\includegraphics[width=1.1\textwidth]{../../version29-with-ocean-masking-lsm/output_2024/track_6/correlation_analysis/relation/track_6_SWVL1-HF_relation.png} \\
			\includegraphics[width=1.1\textwidth]{../../version29-with-ocean-masking-lsm/output_2024/track_6/correlation_analysis/relation/track_6_TP-HF_relation.png}
		\end{center}

		\column{0.50\textwidth}
		\vspace{-0.4cm}
		\begin{center}
			\includegraphics[width=1.1\textwidth]{../../version29-with-ocean-masking-lsm/output_2024/track_6/correlation_analysis/relation/track_6_TP-SWVL1_relation.png}
		\end{center}
	\end{columns}
\end{frame}

% -------------------------------------------------------------------
\subsubsection{Track 6 Hovm\"oller}
\includepdf[
	pages=-,
	scale=1,
	frame=true,
	pagecommand={\stepcounter{framenumber}},
	offset=0mm 0mm
]{../slides/hov-6.pdf}

% -------------------------------------------------------------------
\subsubsection{Track 6 SRTM Elevation Map}
\setLayout{plain}
\begin{frame}{Track 6 SRTM Elevation Map}
	\vspace{-0.4cm}
	\begin{center}
		\includegraphics[width=0.6\textwidth]{../../version29-with-ocean-masking-lsm/output_2024/track_6/elevation_analysis/MD_track_6_elevation_mosaic.png}
	\end{center}
\end{frame}

% -------------------------------------------------------------------
\subsubsection{Track 6 Translation Speed}
\begin{frame}{Track 6 Translation Speed of MD Center}
	\vspace{-0.3cm}
	\begin{center}
		\includegraphics[width=0.8\textwidth]{../../../version29-with-ocean-masking-lsm/output_2024/track_6/translational-speed/track_6_speed_ms.png}

		\begin{block}{\footnotesize Interpretation}
			\scriptsize Track originate over land, shows a steady moderate translation speed that gradually increases as the system moves north-westward across the Indo-Gangetic Plain. The initial deceleration phase over Eastern India allows sufficient precipitation to recharge the already-wet clay soils (\texttt{SWVL1} $\sim 0.40$ m$^3$/m$^3$). The subsequent speed increase due to the soil moisture - precipitation feedback, allowing the system to sustain itself for 117 hours.
		\end{block}
	\end{center}
\end{frame}

% -------------------------------------------------------------------
\section{Observations \& Analysis}
\setLayout{mainpoint}
\begin{frame}{}
	\frametitle{OBSERVATIONS \& ANALYSIS}
\end{frame}

% -------------------------------------------------------------------
\subsection{Diurnal Effects}
\setLayout{vertical}
\begin{frame}{Diurnal Effects in Hovm\"oller Diagrams}
	Strong 24-hour periodic signals appear \textit{exclusively} after land contact.

	\vspace{0.2cm}
	\footnotesize
	\textbf{Variables showing clear diurnal wave patterns:}
	\begin{itemize}
		\item \textbf{Sensible Heat Flux (\texttt{SHF}):} Daytime peaks (\texttt{solar heating}) with amplitude $\sim$56 W/m$^2$ over land vs.\ $\sim$14 W/m$^2$ over ocean.
		\item \textbf{Latent Heat Flux (\texttt{LHF}):} Phase-shifted relative to SHF; evaporation peaks in afternoon but extends through the night (\texttt{soil moisture ``battery''}).
		\item \textbf{Soil Temperature (\texttt{STL1}):} Clear 2–4 K amplitude diurnal cycle.
		\item \textbf{2m Air Temperature (\texttt{T2M}) \& 2m Dew point Temperature (\texttt{D2M}):} Damped oscillation following STL1.
	\end{itemize}

	\vspace{0.2cm}
	\begin{block}{\small Why Land Shows Oscillations but Ocean Does Not}
		\footnotesize
		\begin{center}
			\begin{tabular}{lcc}
				\toprule
				\rowcolor{MyBlue!15}
				Property        & Ocean                                 & Land                \\
				\midrule
				SHF amplitude   & 5–20 W/m$^2$                          & 40–100 W/m$^2$      \\
				Thermal inertia & Very high (20–100m ocean mixed layer) & Low (0–7 cm soil)   \\
				Flux uniformity & High (homogeneous)                    & Low (heterogeneous) \\
				\bottomrule
			\end{tabular}
		\end{center}
	\end{block}

	\begin{alertblock}{\small LHF Phase Shift - Brown Ocean Signature}
		\footnotesize
		Over land, LHF is regulated by \textbf{SWVL1}.
		Maintained high evaporation outside peak heating hours -
		sustaining convection overnight - is the primary signature
		of the Brown Ocean Effect.
	\end{alertblock}
\end{frame}

% -------------------------------------------------------------------
\subsection{Elevation \& Soil Type}
\setLayout{plain}
\begin{frame}{Terrain Control - Elevation Map}
	\begin{columns}[T]
		\column{0.49\textwidth}
		\vspace{-0.5cm}
		\begin{figure}
			\includegraphics[width=0.9\linewidth]{../../version29-with-ocean-masking-lsm/output_2024/global_elevation_analysis/elevation_analysis/MD_all_tracks_elevation.png}
		\end{figure}

		\begin{itemize}\footnotesize
			\item Eastern Ghats: 600–1400 m (Tracks 1, 4).
			\item Indo-Gangetic Plain: $<$400 m (Tracks 2, 3, 5, 6).
		\end{itemize}

		\column{0.49\textwidth}
		\vspace{-0.5cm}
		\begin{figure}
			\includegraphics[width=0.8\linewidth]{../images/Elevation-map-of-India.jpg}
		\end{figure}
		{\tiny \textbf{Fig :} Elevation map of India. Source: Wikimedia -- India topo.}
	\end{columns}

	\vspace{0.2cm}
	\footnotesize \textbf{Orographic disrup feedback mechanism:} Steep slopes ($>$8$^\circ$) increase drainage (down-flow) vs.\ flat terrain $\to$ no infiltration $\to$ soil stays dry $\to$ LHF collapses $\to$ MD decays. \\
	\footnotesize \textit{Tracks over plains survive 2 times longer (mean 135h) than Eastern Ghats tracks (mean 77h).}
\end{frame}

% -------------------------------------------------------------------
\setLayout{plain}
\begin{frame}{Soil Type Control - Plains vs.\ Coastal and Eastern Ghats Soils}
	\begin{columns}[T]
		\column{0.49\textwidth}
		\vspace{-0.4cm}
		\begin{figure}
			\includegraphics[width=0.84\linewidth]{../../version29-with-ocean-masking-lsm/output_2024/MD_track_2024.png}
		\end{figure}

		\column{0.49\textwidth}
		\vspace{-0.4cm}
		\begin{figure}
			\includegraphics[width=0.75\linewidth]{../images/Soiltype-map-of-India.png}
		\end{figure}
		{\tiny \textbf{Fig :} Soil type map of India. Source: European Soil Data Centre (ESDAC).}
	\end{columns}

	\vspace{0.1cm}
	\footnotesize
	\begin{tabular}{lcccc}
		\toprule
		\rowcolor{MyBlue!15}
		\textbf{Region}      & \textbf{Soil Type}         & \textbf{Tracks} & \textbf{Mean Duration} \\
		\midrule
		Eastern coastal belt & Sandy-loam (low retention) & T1, T4          & 77 h                   \\
		Indo-Gangetic Plain  & Clay-rich (high retention) & T2, T3, T5, T6  & 134 h                  \\
		\bottomrule
	\end{tabular}
\end{frame}

% -------------------------------------------------------------------
\section{Conclusion}
\setLayout{mainpoint}
\begin{frame}{}
	\frametitle{CONCLUSION}
\end{frame}

% -------------------------------------------------------------------
\setLayout{plain}
\begin{frame}{Conclusion - Drivers of MD Inland Longevity}
	\vspace{-0.3cm}
	\begin{columns}[T]
		\column{0.49\textwidth}

		\begin{block}{\small \faIcon{tint} Soil Moisture (SWVL1)}
			\footnotesize
			\begin{itemize}
				\item High soil moisture ($>$0.32 m$^3$/m$^3$) over clay-rich IGP soils drives sustained evapotranspiration.
				\item Long-lived tracks (T2, T3, T6): SWVL1 $\sim$0.40–0.45 m$^3$/m$^3$ - maintains moisture feedback to MD core.
				\item Short-lived tracks (T1, T4): SWVL1 $<$0.25 m$^3$/m$^3$ - insufficient to sustain feedback.
				\item \textbf{Soil moisture is the primary energy reservoir inland}.
			\end{itemize}
		\end{block}

		\vspace{0.2cm}
		\begin{block}{\small \faIcon{cloud-rain} Total Precipitation (TP)}
			\footnotesize
			\begin{itemize}
				\item TP after landfall recharges SWVL1, closing the feedback loop.
				\item Positive TP–SWVL1 coupling sustained long-lived tracks.
				\item Without recharge, soil dries $\to$ LHF collapses $\to$ MD decays.
				\item \textbf{TP is the recharge mechanism}.
			\end{itemize}
		\end{block}

		\column{0.51\textwidth}
		\begin{block}{\small \faIcon{fire} Heat Flux (SHF \& LHF)}
			\footnotesize
			\begin{itemize}
				\item Land SHF amplitude stronger than ocean ($\sim$78 W/m$^2$ vs.\ $\sim$14 W/m$^2$)
				\item Diurnal SHF warms the boundary layer $\to$ maintains low-level convergence.
				\item LHF extends evaporation overnight preventing vortex core drying.
				\item \textbf{SHF + LHF diurnal cycle sustaining Brown Ocean}
			\end{itemize}
		\end{block}

		\vspace{0.1cm}
		\begin{block}{\small \faIcon{mountain} Elevation (Terrain)}
			\footnotesize
			\begin{itemize}
				\item Eastern Ghats ($>$900 m, slopes $>$8$^\circ$): rapid runoff prevents infiltration $\to$ feedback broken (Tracks 1, 4: $\sim$77 h).
				\item Indo-Gangetic Plain ($<$400 m): flat terrain promotes infiltration $\to$ \texttt{SWVL1} sustained (Tracks 2, 3, 5, 6: $\sim$134 h).
				\item \textbf{Plains enable while high elevation (Eastern Ghats) disable the Brown Ocean Effect}.
			\end{itemize}
		\end{block}
	\end{columns}

	\vspace{0.1cm}
	\begin{alertblock}{}
		\centering \footnotesize
		\textbf{MD inland longevity} $=$ \textit{Wet soil} (SWVL1) $+$ \textit{Precipitation recharge} (TP) $+$ \textit{Diurnal heat flux pumping} (SHF/LHF) $+$ \textit{Flat terrain} (elevation)
	\end{alertblock}
\end{frame}
\setLayout{horizontal}
\begin{frame}{Limitations \& Future Work}
	\vspace{-0.5cm}
	\begin{block}{\small Current Limitations}
		\begin{itemize}\footnotesize
			\item \textbf{Sample size:} 6 tracks (1 season, 2024). Statistically limited; multi-year analysis needed.
			\item \textbf{ERA5 resolution:} 28 km smooths soil heterogeneity; coastal SWVL1 may be underestimated.
			      % \item \textbf{Causality:} TP–SWVL1 correlation cannot confirm direction (Granger test failed due to short track length 100–140 h).
			\item \textbf{Harmonized world soil database (HWSD) integration:} Soil content (Clay, sand, etc.) estimated for along each track region; direct soil texture mapping.
			\item 850 hPa vorticity analysis alongside 10m and 100m wind component fields.
			\item \textbf{Sea Surface Temperature (\texttt{SST}) effects:} Late-season track differences may partly reflect SST gradient changes, not only soil.
		\end{itemize}
	\end{block}
	\vfill
\end{frame}

% -------------------------------------------------------------------
\section{References}
\setLayout{vertical}
\begin{frame}{References}
	\vspace{-0.2cm}
	\footnotesize
	\textbf{ERA5 Data Source}
	\begin{itemize}\scriptsize
		\item Hersbach, H., Bell, B., Berrisford, P., et al. (2020). The ERA5 global reanalysis. \textit{Quarterly Journal of the Royal Meteorological Society}, 146(730), 1999–2049.
	\end{itemize}

	\vspace{-0.1cm}
	\footnotesize
	\textbf{Soil Moisture-Precipitation Relationship}
	\begin{itemize}\scriptsize
		\item Jose, V. \& Chandrasekar, A. (2025). Long-term relationship between soil moisture and precipitation over India: An analysis using event coincidence analysis and PCMCI method. \textit{Theoretical and Applied Climatology}, 156. doi:10.1007/s00704-025-05524-w
	\end{itemize}

	\vspace{-0.1cm}
	\footnotesize
	\textbf{Brown Ocean Effect}
	\begin{itemize}\scriptsize
		\item Vijay, D. V. S. \& Chandrasekar, A. (2024). A study on the possible Brown Ocean Effect: Impacts of an antecedent tropical cyclone on the rapid intensification of the 1999 Odisha Supercyclone. \textit{Journal of Geophysical Research: Atmospheres}, 129(1), e2023JD039409. doi:10.1029/2023JD039409
		\item Nair, U. S., Rappin, E., Foshee, E., et al. (2019). Influence of Land Cover and Soil Moisture based Brown Ocean Effect on an Extreme Rainfall Event from a Louisiana Gulf Coast Tropical System. \textit{Scientific Reports}, 9, 17136. doi:10.1038/s41598-019-53031-6
		\item Andersen, T. K. \& Shepherd, J. M. (2013). A global climatology of inland tropical cyclone maintenance or intensification. \textit{Journal of Climate}, 26(17), 6629–6642.
		\item Andersen, T. K. \& Marshall Shepherd, J. (2014). Insights into the Brown Ocean Effect: A Case Study of Tropical Storm Erin (2007). \textit{International Journal of Climatology}, 34(5), 1607–1615. doi:10.1002/joc.3784
		\item Wakefield, R. A., Rapp, A. D., Shepherd, J. M., \& Santanello Jr, J. A. (2021). The inland maintenance of Tropical Storm Bill (2015). Part I: Contributions of the Brown Ocean Effect. \textit{Monthly Weather Review}, 149(9), 3083–3102.
	\end{itemize}
\end{frame}

\begin{frame}{References}
	\vspace{-0.2cm}
	\footnotesize
	\textbf{Soil Moisture Initialization \& MD Prediction}
	\begin{itemize}\scriptsize
		\item Osuri, K. K., Nadimpalli, R., Mohanty, U. C., \& Niyogi, D. (2020). Improved simulation of monsoon depressions from soil moisture initialization. \textit{Journal of Geophysical Research: Atmospheres}, 125(18), e2020JD032400.
		\item Lodh, A., Routray, A., Dutta, D., George, J. P., \& Mitra, A. K. (2022). Improving the prediction of monsoon depressions by assimilating ASCAT soil moisture in NCUM-R modeling system. \textit{Atmospheric Research}, 272, 106130. doi:10.1016/j.atmosres.2022.106130
	\end{itemize}

	\vspace{-0.1cm}
	\footnotesize
	\textbf{Soil Moisture Perturbations \& Vortex Intensity}
	\begin{itemize}\scriptsize
		\item Hunt, K. M. \& Turner, A. G. (2017). The effect of soil moisture perturbations on Indian monsoon depressions in a numerical weather prediction model. \textit{Journal of Climate}, 30(21), 8803–8823.
	\end{itemize}

	\vspace{-0.1cm}
	\footnotesize
	\textbf{Land Surface Feedback \& Post-Landfall Structure}
	\begin{itemize}\scriptsize
		\item Chang, H.-Y., Niyogi, D., Mohanty, U. C., Chen, F., \& Leung, L. R. (2009). Possible relation between land surface feedback and the post-landfall structure of monsoon depressions. \textit{Geophysical Research Letters}, 36(15), L15821.
		\item Nair, U. S., Rapp, A. D., Foshee, E., et al. (2011). Influence of land surface processes on the intensification of land-falling tropical cyclones. \textit{Monthly Weather Review}, 139(5), 1409–1428.
	\end{itemize}

	\vspace{-0.1cm}
	\footnotesize
	\textbf{Soil Hydraulic Properties}
	\begin{itemize}\scriptsize
		\item Wood, E. F. \& Roundy, R. (2005). Using soil hydraulic properties to predict soil water dynamics. \textit{Advances in Water Resources}, 28(7), 757–770.
		\item Clapp, R. B. \& Hornberger, G. M. (1978). Empirical equations for some soil hydraulic properties. \textit{Water Resources Research}, 14(4), 601–604.
	\end{itemize}
\end{frame}

\begin{frame}{References}
	\vspace{-0.2cm}
	\footnotesize
	\textbf{Land Use/Land Cover Impact}
	\begin{itemize}\scriptsize
		\item Ankur, K. \& Osuri, K. K. (2022). An assessment study on the impact of land use land cover on monsoon depressions over India. \textit{Atmospheric Research}, 279, 106375. doi:10.1016/j.atmosres.2022.106375
	\end{itemize}

	\vspace{-0.1cm}
	\footnotesize
	\textbf{Pre-Storm Ground Wetness \& Inland Penetration}
	\begin{itemize}\scriptsize
		\item Kishtawal, C. M., Niyogi, D., Tewari, M., Pielke Sr, R. A., \& Beltrán-Przekurat, A. (2013). Enhancement of inland penetration of monsoon depressions in the Bay of Bengal due to prestorm ground wetness. \textit{Water Resources Research}, 49(6), 3573–3589.
		\item Kishtawal, C. M., Niyogi, D., Tewari, M., Pielke Sr, R. A., \& Shepherd, J. M. (2013). Tropical cyclone intensification over land after landfall. \textit{Nature}, 493, 209–212.
	\end{itemize}

	\vspace{-0.1cm}
	\footnotesize
	\textbf{Soil Moisture Memory \& Climate Interactions}
	\begin{itemize}\scriptsize
		\item Seneviratne, S. I., Corti, T., Davin, E. L., et al. (2010). Investigating soil moisture–climate interactions in a changing climate: A review. \textit{Earth-Science Reviews}, 99(3-4), 125–161.
	\end{itemize}

	\vspace{-0.1cm}
	\footnotesize
	\textbf{Monsoon Depression Case Studies}
	\begin{itemize}\scriptsize
		\item Mishra, S. K. (2021). Development of a Monsoon Depression and Its Interaction with the Large-Scale Background: A Case Study. \textit{Monthly Weather Review}, 149(11), 3627–3644.
	\end{itemize}

	\vspace{-0.1cm}
	\footnotesize
	\textbf{Orography \& Cloud Dynamics}
	\begin{itemize}\scriptsize
		\item Houze Jr, R. A. (1994). \textit{Cloud Dynamics} (2nd ed.). Academic Press.
		\item Tuttle, S. \& Davis, C. A. (2006). Lessons from 1997–2005 Kansas-Nebraska high plains drought. \textit{Journal of Applied Meteorology and Climatology}, 45(2), 221–242.
		\item Kirshbaum, D. J. (2018). Environmental controls on orographic cloud evolution. \textit{Journal of Geophysical Research: Atmospheres}, 123(14), 7421–7434.
	\end{itemize}
\end{frame}

% -------------------------------------------------------------------
\section{Thank You}
\setLayout{blank}
\begin{frame}
	\centering
	\vspace{3.5cm}
	\textbf{\Huge Thank You} \\
	\vspace{1cm}
	\textbf{Questions and Suggestions}
\end{frame}

\end{document}
